\SourceLeaf{13}
\begin{center}\Large EIN NEUER BEWEIS\end{center}

\begin{center}\large EINES ARITHMETISCHEN SATZES.\end{center}

\BookRule{35mm}

\Article{1}

Fragen der höheren Arithmetik zeigen sehr häufig eine eigentümliche Erscheinung, die in der Analysis viel seltener vorkommt und ihren Reiz erheblich steigert. Bei analytischen Untersuchungen können wir meist erst dann zu neuen Wahrheiten gelangen, wenn wir die Prinzipien, auf denen sie beruhen und die uns gleichsam den Weg zu ihnen öffnen müssen, zuvor gründlich beherrschen. In der Arithmetik dagegen treten sehr häufig durch Induktion und eine Art unerwarteten Glücks neue Wahrheiten von größter Eleganz hervor, deren Beweise so tief verborgen und in so große Dunkelheit gehüllt sind, dass sie sich allen Versuchen entziehen und selbst den scharfsinnigsten Untersuchungen den Zugang verwehren. Außerdem besteht zwischen arithmetischen Wahrheiten, die auf den ersten Blick völlig verschiedenartig erscheinen, ein so großer und wunderbarer Zusammenhang, dass wir nicht selten, während wir ganz anderes suchen, schließlich auf das Glücklichste zu dem so sehr ersehnten und zuvor in langen Überlegungen vergeblich gesuchten Beweis gelangen, und zwar auf einem ganz anderen Weg als dem erwarteten. Meist sind solche Wahrheiten aber von der Art, dass man auf mehreren sehr verschiedenen Wegen zu ihnen gelangen kann, und nicht immer sind die Wege, die sich zuerst anbieten, die kürzesten. Es muss daher gewiss hoch geschätzt werden, wenn es nach langer vergeblicher Untersuchung einer solchen Wahrheit, die dann zwar bewiesen wurde, aber auf ziemlich verborgenen Umwegen, schließlich gelingt, den einfachsten und eigentlichen Weg zu entdecken.

\Article{2}

Unter den im vorigen Artikel besprochenen Fragen nimmt der Satz einen hervorragenden Platz ein, der fast die gesamte Theorie der quadratischen Reste enthält und in den \emph{Untersuchungen über höhere Arithmetik} (Abschnitt IV) mit dem Namen \emph{Fundamentalsatz} ausgezeichnet \SourceLeaf{14}ist. Als \emph{erster} Entdecker dieses höchst eleganten Satzes ist ohne Zweifel der berühmte \textsc{Legendre} anzusehen, nachdem die großen Geometer \textsc{Euler} und \textsc{Lagrange} schon lange zuvor mehrere seiner besonderen Fälle durch Induktion entdeckt hatten. Ich halte mich hier nicht damit auf, die Beweisversuche dieser Männer aufzuzählen; wer möchte, möge das eben erwähnte Werk nachlesen. Zur Bestätigung des im vorigen Artikel Gesagten sei mir nur gestattet hinzuzufügen, was meine eigenen Versuche betrifft. Ich war im Jahr 1795 selbstständig auf den Satz gestoßen, als ich von allem, was in der höheren Arithmetik schon ausgearbeitet worden war, noch gar nichts wusste und von allen Hilfsmitteln der Literatur völlig abgeschnitten war. Doch ein ganzes Jahr lang quälte er mich und entzog sich meinen angestrengtesten Bemühungen, bis ich schließlich den im vierten Abschnitt jenes Werkes gegebenen Beweis fand. Später boten sich mir drei weitere an, die auf ganz verschiedenen Prinzipien beruhen; einen davon habe ich im fünften Abschnitt gegeben, die übrigen, die ihm an Eleganz nicht nachstehen, werde ich bei anderer Gelegenheit veröffentlichen. Doch alle diese Beweise scheinen zwar hinsichtlich ihrer Strenge nichts zu wünschen übrigzulassen, sind aber aus zu verschiedenartigen Prinzipien hergeleitet, vielleicht mit Ausnahme des ersten, der jedoch mit recht mühsamen Überlegungen vorgeht und durch ziemlich weitläufige Rechnungen belastet ist. Ich zögere daher nicht auszusprechen, dass bisher ein \emph{eigentlicher} Beweis gefehlt hat. Nun mögen die Sachkundigen beurteilen, ob der Beweis, den ich vor Kurzem entdeckt habe und den die folgenden Seiten enthalten, diesen Namen verdient.

\Article{3}

\textsc{Satz.} \emph{Es sei $p$ eine positive Primzahl; $k$ eine beliebige, durch $p$ nicht} \emph{teilbare ganze Zahl; es seien}

\[
\begin{aligned}
A &\quad \text{\emph{Gesamtheit der Zahlen}}\quad 1,\ 2,\ 3\ \dotsfour\ \tfrac12(p-1)\\
B &\quad \text{\emph{Gesamtheit dieser Zahlen}}\quad \tfrac12(p+1),\ \tfrac12(p+3),\ \tfrac12(p+5)\ \dotsfour\ p-1
\end{aligned}
\]

\emph{Man bilde die Produkte aus $k$ und den einzelnen Zahlen in $A$ und nehme ihre kleinsten positiven Reste modulo $p$, die offenbar alle verschieden sind und teils zu $A$, teils zu} \emph{$B$ gehören. Gehören nun zu $B$ insgesamt $\mu$ Reste, so ist $k$ ein quadratischer Rest oder Nichtrest von $p$, je nachdem $\mu$ gerade oder ungerade ist.}

\emph{Beweis.} Die zu $A$ gehörenden Reste seien $a,\ a'\ a''\ \dotsfour$, die übrigen, zu $B$ gehörenden $b,b',b''\dotsfour$. Offenbar sind die Ergänzungen der letzteren, $p-b,p-b',p-b''\dotsfour$, sämtlich von den Zahlen $a,\ a',\ a''\ \dotsfour$ verschieden, füllen aber mit diesen zusammen die Gesamtheit\SourceLeaf{15} $A$ aus. Wir haben also

\[
\begin{aligned}
1\pdot2\pdot3\ \dotsfour\ \tfrac12(p-1)=aa' a''\ \dotsfour\ (p-b)(p-b')(p-b'')\ \dotsfour
\end{aligned}
\]

Das letztere Produkt wird aber offenbar

\[
\begin{aligned}
&\equiv(-1)^\mu aa' a''\ \dotsfour\ bb' b''\ \dotsfour\ \equiv(-1)^\mu k\pdot2k\pdot3k\ \dotsfour\ \tfrac12(p-1)k\\
&\equiv(-1)^\mu k^{\frac12(p-1)}\,1\pdot2\pdot3\ \dotsfour\ \tfrac12(p-1) (\mathrm{mod.}p)
\end{aligned}
\]

\Needspace{7\baselineskip}Daher ist

\[
\begin{aligned}
1\equiv(-1)^\mu k^{\frac12(p-1)}
\end{aligned}
\]

oder $k^{\frac12(p-1)}\equiv\pm1$, je nachdem $\mu$ gerade oder ungerade ist; daraus folgt unser Satz unmittelbar.

\Article{4}

Durch die Einführung einiger geeigneter Bezeichnungen können wir die folgenden Überlegungen erheblich verkürzen. Das Zeichen $(k,p)$ soll uns daher die Anzahl derjenigen Produkte unter

\[
\begin{aligned}
k,\ 2k,\ 3k\ \dotsfour\ \tfrac12(p-1)k,
\end{aligned}
\]

bezeichnen, deren kleinste positive Reste modulo $p$ dessen Hälfte übersteigen. Ist ferner $x$ eine beliebige nicht ganze Größe, so bezeichnen wir mit dem Zeichen $[x]$ die unmittelbar unter $x$ liegende ganze Zahl, sodass $x-[x]$ stets eine positive Größe zwischen den Grenzen $0$ und $1$ ist. Die folgenden Beziehungen ergeben sich nun mit geringer Mühe:

I. $[x]+[-x]=-1.$

II. $[x]+h=[x+h]$, wenn $h$ eine ganze Zahl ist.

III. $[x]+[h-x]=h-1.$

IV. Ist $x-[x]$ ein Bruch kleiner als $\tfrac12$, so ist $[2x]-2[x]=0$; ist aber $x-[x]$ größer als $\tfrac12$, so ist $[2x]-2[x]=1.$

V. Liegt also der kleinste positive Rest der ganzen Zahl $h$ modulo $p$ unter $\tfrac12p$, so ist $\gbr{\frac{2h}{p}}-2\gbr{\frac{h}{p}}=0$; liegt jener Rest aber über $\tfrac12p$, so ist \mbox{$\gbr{\frac{2h}{p}}-2\gbr{\frac{h}{p}}=1.$}

\SourceLeaf{16}
VI. Daraus folgt sogleich $(k,p)=$

\[
\begin{aligned}
&\gbr{\frac{2k}{p}}+\gbr{\frac{4k}{p}}+\gbr{\frac{6k}{p}}\dotsfour+\gbr{\frac{(p-1)k}{p}}\\
&-2\gbr{\frac{k}{p}}-2\gbr{\frac{2k}{p}}-2\gbr{\frac{3k}{p}}\dotsfour-2\gbr{\frac{\frac12(p-1)k}{p}}.
\end{aligned}
\]

VII. Aus VI und I ergibt sich ohne Mühe

\[
\begin{aligned}
(k,p)+(-k,p)=\tfrac12(p-1)
\end{aligned}
\]

Daraus folgt, dass $-k$ dieselbe oder die entgegengesetzte Beziehung zu $p$ hat (insofern es dessen quadratischer Rest oder Nichtrest ist) wie $+k$, je nachdem $p$ die Form $4n+1$ oder die Form $4n+3$ hat. Im ersten Fall ist $-1$ offenbar ein Rest, im zweiten ein Nichtrest von $p$.

\Needspace{10\baselineskip}VIII. Die in VI gegebene Formel werden wir folgendermaßen umformen. Nach III ist

\[
\begin{aligned}
\gbr{\frac{(p-1)k}{p}}&=k-1-\gbr{\frac{k}{p}},\\
\gbr{\frac{(p-3)k}{p}}&=k-1-\gbr{\frac{3k}{p}},\\
\gbr{\frac{(p-5)k}{p}}&=k-1-\gbr{\frac{5k}{p}}\dotsfour
\end{aligned}
\]

Wenden wir diese Ersetzungen auf die letzten $\frac{p\mp1}{4}$ Glieder der oberen Reihe in jenem Ausdruck an, so erhalten wir

\Needspace{8\baselineskip}\emph{erstens}, wenn $p$ die Form $4n+1$ hat,

\[
\begin{aligned}
(k,p)&=\tfrac14(k-1)(p-1)\\
& -2\left\{\gbr{\frac{k}{p}}+\gbr{\frac{3k}{p}}+\gbr{\frac{5k}{p}}\dotsfour+\gbr{\frac{\frac12(p-3)k}{p}}\right\}\\
&\qquad-\left\{\gbr{\frac{k}{p}}+\gbr{\frac{2k}{p}}+\gbr{\frac{3k}{p}}\dotsfour+\gbr{\frac{\frac12(p-1)k}{p}}\right\}
\end{aligned}
\]

\Needspace{8\baselineskip}\emph{zweitens}, wenn $p$ die Form $4n+3$ hat,

\[
\begin{aligned}
(k,p)&=\tfrac14(k-1)(p+1)\\
& -2\left\{\gbr{\frac{k}{p}}+\gbr{\frac{3k}{p}}+\gbr{\frac{5k}{p}}\dotsfour+\gbr{\frac{\frac12(p-1)k}{p}}\right\}\\
&\qquad-\left\{\gbr{\frac{k}{p}}+\gbr{\frac{2k}{p}}+\gbr{\frac{3k}{p}}\dotsfour+\gbr{\frac{\frac12(p-1)k}{p}}\right\}
\end{aligned}
\]

IX. Für den besonderen Fall $k=+2$ folgt aus den eben gegebenen Formeln $(2,p)=\tfrac14(p\mp1)$, wobei das obere oder untere Zeichen zu nehmen ist, je nachdem $p$ die Form $4n+1$ oder $4n+3$ hat. Somit ist $(2,p)$ gerade, also $2Rp$, wenn $p$ die Form $8n+1$ oder $8n+7$ hat; dagegen ist $(2,p)$ ungerade und $2Np$, wenn $p$ die Form $8n+3$ oder $8n+5$ hat.

\SourceLeaf{17}
\Needspace{14\baselineskip}\Article{5}

\textsc{Satz.} \emph{Es sei $x$ eine positive nicht ganze Größe, unter deren Vielfachen $x$, $2x$,} \emph{$3x\dotsfour$ bis $nx$ keines eine ganze Zahl ist; man setze $[nx]=h$, woraus leicht folgt,} \emph{dass auch unter den Vielfachen der reziproken Größe $\frac1x$, $\frac2x$, $\frac3x\dotsfive$ bis $\frac hx$ keine ganze Zahl} \emph{vorkommt. Dann behaupte ich, es sei}

\[
\begin{aligned}
\left.\begin{aligned}
&[x]+[2x]+[3x]\dotstwo+[nx]\\
&+\gbr{\frac1x}+\gbr{\frac2x}+\gbr{\frac3x}\dotsfour+\gbr{\frac hx}
\end{aligned}\right\}=nh.
\end{aligned}
\]

\emph{Beweis.} In der Reihe $[x]+[2x]+[3x]\dotsfour+[nx]$, die wir $=\Omega$ setzen, sind die ersten Glieder bis einschließlich des Gliedes an der Stelle $\gbr{\frac1x}$ offenbar alle $=0$; die folgenden bis zum Glied an der Stelle $\gbr{\frac2x}$ sämtlich $=1$; die folgenden bis zum Glied an der Stelle $\gbr{\frac3x}$ sämtlich $=2$, und so weiter. Daher ist

\[
\begin{aligned}
\left.\begin{aligned}
\Omega={}&0\times\gbr{\frac1x}\\
&+1\times\left\{\gbr{\frac2x}-\gbr{\frac1x}\right\}\\
&+2\times\left\{\gbr{\frac3x}-\gbr{\frac2x}\right\}\\
&+3\times\left\{\gbr{\frac4x}-\gbr{\frac3x}\right\}\\
&\hspace{12mm}\text{usw.}\\
&+(h-1)\left\{\gbr{\frac hx}-\gbr{\frac{h-1}{x}}\right\}\\
&+h\left\{n-\gbr{\frac hx}\right\}
\end{aligned}\right\}=hn-\gbr{\frac1x}-\gbr{\frac2x}-\gbr{\frac3x}\dotsfour-\gbr{\frac hx}
\end{aligned}
\]

Q. E. D.

\Needspace{12\baselineskip}\Article{6}

\textsc{Satz.} \emph{Bezeichnen $k,p$ beliebige positive ungerade, zueinander teilerfremde Zahlen, so ist}

\[
\begin{aligned}
\left.\begin{aligned}
&\gbr{\frac kp}+\gbr{\frac{2k}{p}}+\gbr{\frac{3k}{p}}\dotsfour+\gbr{\frac{\frac12(p-1)k}{p}}\\
&+\gbr{\frac pk}+\gbr{\frac{2p}{k}}+\gbr{\frac{3p}{k}}\dotsfour+\gbr{\frac{\frac12(k-1)p}{k}}
\end{aligned}\right\}=\tfrac14(k-1)(p-1).
\end{aligned}
\]

\emph{Beweis.} Nehmen wir, was erlaubt ist, $k<p$ an, so ist $\frac{\frac12(p-1)k}{p}$ kleiner als $\frac12 k$, aber größer als $\frac12(k-1)$, also $\gbr{\frac{\frac12(p-1)k}{p}}=\frac12(k-1)$. Daraus ergibt sich offenbar der vorliegende Satz unmittelbar aus dem vorhergehenden, wenn man dort $\frac kp=x$, $\frac12(p-1)=n$, also $\frac12(k-1)=h$ setzt.

\SourceLeaf{18}
\Needspace{8\baselineskip}Übrigens lässt sich auf ähnliche Weise beweisen, dass, wenn $k$ eine zu $p$ teilerfremde \emph{gerade} Zahl ist, gilt:

\[
\begin{aligned}
\left.\begin{aligned}
&\gbr{\frac kp}+\gbr{\frac{2k}{p}}+\gbr{\frac{3k}{p}}\dotsfour+\gbr{\frac{\frac12(p-1)k}{p}}\\
&+\gbr{\frac pk}+\gbr{\frac{2p}{k}}+\gbr{\frac{3p}{k}}\dotsfour+\gbr{\frac{\frac12 kp}{k}}
\end{aligned}\right\}=\tfrac14 k(p-1)
\end{aligned}
\]

Doch halten wir uns bei diesem für unser Vorhaben nicht notwendigen Satz nicht auf.

\clearpage\Article{7}

Aus der Verbindung des vorhergehenden Satzes mit der Aussage VIII des Artikels 4 folgt nun unmittelbar der Fundamentalsatz. Bezeichnen nämlich $k,p$ beliebige verschiedene positive Primzahlen, und setzen wir

\[
\begin{aligned}
(k,p)+\gbr{\frac kp}+\gbr{\frac{2k}{p}}+\gbr{\frac{3k}{p}}\dotsfour+\gbr{\frac{\frac12(p-1)k}{p}}&=L\\
(p,k)+\gbr{\frac pk}+\gbr{\frac{2p}{k}}+\gbr{\frac{3p}{k}}\dotsfour+\gbr{\frac{\frac12(k-1)p}{k}}&=M
\end{aligned}
\]

so ist nach VIII des Artikels 4 klar, dass $L$ und $M$ stets gerade Zahlen sind. Nach dem Satz des Artikels 6 ist aber

\[
\begin{aligned}
L+M=(k,p)+(p,k)+\tfrac14(k-1)(p-1)
\end{aligned}
\]

Wann immer also $\frac14(k-1)(p-1)$ gerade ist, was geschieht, wenn für beide $k,p$ oder wenigstens für eine von ihnen die Form $4n+1$ vorliegt, müssen $(k,p)$ und $(p,k)$ notwendig entweder beide gerade oder beide ungerade sein. Ist hingegen $\frac14(k-1)(p-1)$ ungerade, was geschieht, wenn beide $k,p$ die Form $4n+3$ haben, so muss notwendig eine der Zahlen $(k,p)$, $(p,k)$ gerade, die andere ungerade sein. Im ersten Fall sind daher die Beziehung von $k$ zu $p$ und die Beziehung von $p$ zu $k$ (insofern eine der beiden ein Rest oder Nichtrest der anderen ist) gleich, im zweiten entgegengesetzt.

Q. E. D.

\BookRule{35mm}
